% methods

\section{Study Design}\label{sec:design}

\begin{figure}[htbp]
  \centering
  % Replace the \fbox with:
  % \includegraphics[width=\textwidth]{Figures/pipeline.pdf}
  \fbox{\parbox[c][2.5in][c]{0.9\textwidth}{\centering
    \note{Placeholder: modeling framework and pipeline slide}}}
  \caption{Modeling framework: energy-system modeling, emissions processing,
    air-quality modeling, and health-impact assessment.}
  \label{fig:pipeline}
\end{figure}

The study links four modules (Figure~\ref{fig:pipeline}).

\section{Energy-System Modeling}\label{sec:energy}

We use GCAM-KAIST to model different decarbonization pathways (e.g.\ no
policy, linear reduction, early reduction). At a high level, we compare linear
and early-reduction pathway curves with the same endpoints, matching what is
written into law to keep the analysis policy-relevant.

The main points of interest are temporal and regional (provincial) trade-offs,
less so cumulative outcomes.
\note{This differs from the research question in Section~\ref{sec:rq}, which
focuses on cumulative outcomes. Align the two.}

\discuss{Which exact scenarios can we get from GCAM-KAIST?}

\subsection{Facility Retirement Ordering}\label{sec:retirement}

Each reduction pathway likely requires micro-level decisions about which
facilities (e.g.\ power plant units) retire and when. We cannot simply
decrease all units in a sector proportionally: that is unrealistic without
retirements, and emissions are not linear in activity (there may be a
non-zero intercept). We propose a retirement ordering algorithm based on
unit age, emissions, most recent historical activity, and other
characteristics.

\section{Emissions Processing}\label{sec:emissions}

GCAM-KAIST energy outputs are translated into precursor emissions and
downscaled to a gridded level. We apply the same downscaling method as in
\note{cite One Earth paper}: take the national energy structure from
GCAM-KAIST, multiply by emission factors to obtain national emissions, then
downscale using a spatial allocation. We could potentially use provincial
emission factors when downscaling.

Analogous to the use of the U.S.\ National Emissions Inventory (NEI), we
propose using the Clean Air Policy Support System (CAPSS) National Air
Pollutant Emissions Inventory for the spatial allocation mapping from sector
activity to emissions.

\discuss{Resolution of GCAM-KAIST outputs. Should we downscale energy metrics
and then apply emission factors, or downscale processed emissions?}

\section{Air-Quality Modeling}\label{sec:aq}

We use Global InMAP to estimate annual average \PM{} concentrations in Korea.

\discuss{How to treat emissions outside Korea.}
\note{Yehyun: Since Global InMAP is global, doesn't it already account for
transboundary transport, if we also roughly model China?}

Candidate scenario designs for the main analysis:
\begin{itemize}
  \item Korea and other countries (China, India) all reach net zero.
  \item Korea reaches net zero; other countries do not.
  \item A mismatch in net-zero implementation across countries could also be
        interesting to examine.
\end{itemize}

\section{Health-Impact Assessment}\label{sec:hia}

We combine gridded \PM{} exposure with projected population, age structure,
baseline mortality rates, and established concentration--response functions.
